\documentclass[a4paper,12pt]{article}

\usepackage[T1]{fontenc}
\usepackage[italian]{babel}
\usepackage{graphicx}
\usepackage{geometry}
\usepackage{tabularx}
\usepackage[table]{xcolor}
\usepackage[most]{tcolorbox}
\usepackage{fancyhdr}
\usepackage[hidelinks]{hyperref}

\newcommand{\groupmail}{alittlebyte19@gmail.com}
\newcommand{\grouplogo}{../../../assets/logo_alittlebyte.png}
\newcommand{\groupsymbol}{../../../assets/solo_logo.png}

\geometry{
    left=2.5cm,
    right=2.5cm,
    top=2.5cm,
    bottom=2.5cm,
    headheight=331pt,
    includeheadfoot
}

\pagestyle{fancy}
\fancyhf{}

\renewcommand{\headrulewidth}{0pt}
\fancyhead{}

\renewcommand{\footrulewidth}{0.4pt}
\fancyfoot[L]{\includegraphics[height=0.6cm]{\groupsymbol}\hspace{0.45em}aLittleByte}
\fancyfoot[R]{Analisi dei Capitolati}

\fancypagestyle{firstpage}{
    \fancyhf{}
    \renewcommand{\headrulewidth}{0pt}
    \renewcommand{\footrulewidth}{0.4pt}
    \fancyhead[C]{%
        \makebox[\textwidth][c]{%
            \begin{tabular}{c}
                \includegraphics[height=10cm]{\grouplogo}\\[1ex]
                \small\texttt{\groupmail}
            \end{tabular}%
        }
    }
    \fancyfoot[L]{\includegraphics[height=0.6cm]{\groupsymbol}\hspace{0.45em}aLittleByte}
    \fancyfoot[R]{Analisi dei Capitolati}
}

\begin{document}

\thispagestyle{firstpage}

\vspace*{0.5cm}

\begin{center}
    {\LARGE \textbf{Analisi dei Capitolati}}
\end{center}

\vspace{1.5cm}

\begin{center}
    \renewcommand{\arraystretch}{1.3}
    \begin{tcolorbox}[
        width=0.92\textwidth,
        colback=gray!6,
        colframe=gray!45,
        boxrule=0.5pt,
        arc=2mm,
        boxsep=0pt,
        left=8pt, right=8pt, top=8pt, bottom=8pt
    ]
        \begin{tabularx}{\linewidth}{>{\bfseries}p{3.5cm} X}
            Gruppo      & aLittleByte (gruppo 19) \\
            Redattore   & Angelica Gastaldello \\
            Verificatori& \\
            Destinatari & Prof. Tullio Vardanega, Prof. Riccardo Cardin \\
            Data        & \\
        \end{tabularx}
    \end{tcolorbox}
\end{center}

\newpage

\newgeometry{left=2.5cm,right=2.5cm,top=2.5cm,bottom=2.5cm,includefoot}

\begin{center}
    {\Large \textbf{Registro delle Modifiche}}
\end{center}

\vspace{0.35cm}

\renewcommand{\arraystretch}{1.35}
\rowcolors{2}{gray!8}{white}
\begin{center}
    {\footnotesize
    \begin{tabularx}{\textwidth}{|p{1.4cm}|p{1.8cm}|X|p{2.4cm}|p{2.3cm}|p{2.3cm}|}
        \hline
        \rowcolor{gray!20}
        \textbf{Versione} & \textbf{Data} & \textbf{Descrizione} & \textbf{Redattore} & \textbf{Verificatore} & \textbf{Approvatore} \\
        \hline
        0.0.4 & 10/03/2026 & Terminato C1, iniziato C2 & Angelica Gastaldello & - & - \\
        \hline
        0.0.3 & 09/03/2026 & Iniziato C1 & Angelica Gastaldello & - & - \\
        \hline
        0.0.2 & 08/03/2026 & Introduzione e Criteri & Angelica Gastaldello & - & - \\
        \hline
        0.0.1 & 07/03/2026 & Struttura del documento & Angelica Gastaldello & - & - \\
        \hline
    \end{tabularx}
    }
\end{center}

\newpage

\begin{center} 
    {\Large \textbf{Indice}}
\end{center}

\vspace{0.6cm}

\tableofcontents

\newpage

\pagenumbering{arabic}
\setcounter{page}{4}
\fancyfoot[R]{Analisi dei Capitolati\hspace{0.8em}\thepage}

\section{Introduzione}
\subsection{Finalità del documento}

Il presente documento ha lo scopo di illustrare l'analisi e la valutazione dei capitolati d'appalto proposti per l'anno accademico 2025/2026. Al suo interno viene fornita una disamina dettagliata del processo decisionale che ha portato il gruppo aLittleByte alla scelta finale.

Nello specifico, il documento approfondisce le motivazioni tecniche e organizzative alla base della selezione del capitolato C5 - Nexum, proposto dall'azienda Eggon S.r.l. Il progetto prevede lo sviluppo di moduli innovativi di AI Generativa per una piattaforma HR evoluta, integrando tecnologie cloud all'avanguardia. Per garantire la completezza della trattazione, viene inoltre riportata una valutazione comparativa degli altri capitolati disponibili, evidenziando i criteri di confronto utilizzati dal team.

\section{Criteri e metodologia}

La scelta del capitolato da parte del gruppo ha seguito un approccio analitico e strutturato, basato sulla definizione di precisi parametri di valutazione e sull'applicazione di un processo decisionale articolato in più fasi.

\subsection{Fattori considerati nella scelta}

Per valutare e confrontare oggettivamente le diverse proposte aziendali, il team ha individuato i seguenti criteri fondamentali:

\begin{itemize}
    \item \textbf{Complessità e fattibilità:} si è cercato un attento bilanciamento tra il livello di sfida tecnica (in termini di architettura e integrazione di sistemi) e la reale sostenibilità del carico di lavoro. L'obiettivo era garantire un prodotto finale di qualità nel pieno rispetto del monte ore a disposizione.
    \item \textbf{Interesse:} è stato dato un peso rilevante alla curiosità, alla motivazione e al coinvolgimento del gruppo verso il dominio applicativo del progetto, considerando la motivazione condivisa un elemento determinante.
    \item \textbf{Chiarezza delle tecnologie:} si è valutata la trasparenza e la precisione dei requisiti tecnologici richiesti, un aspetto cruciale per stimare preventivamente la curva di apprendimento necessaria e agevolare la futura divisione dei compiti.
    \item \textbf{Disponibilità del proponente:} è stata considerata essenziale la presenza di referenti tecnici aziendali pronti a fornire un supporto continuo, al fine di mitigare i rischi e chiarire tempestivamente eventuali dubbi in fase di sviluppo.
\end{itemize}

\subsection{Processo decisionale}

L'applicazione dei criteri appena descritti è avvenuta attraverso una metodologia progressiva di selezione, che ha progressivamente ristretto il campo delle opzioni:

\begin{itemize}
    \item \textbf{Studio individuale e discussione collettiva:} la fase iniziale ha previsto una lettura e un'analisi autonoma dei singoli capitolati da parte di ogni membro. A questo è seguita una prima valutazione di gruppo, fondamentale per confrontare le opinioni, discutere i punti di forza e di debolezza di ogni proposta e far emergere le preferenze condivise del team.
    \item \textbf{Colloqui con le aziende e decisione finale:} sulla base della scrematura iniziale, il gruppo ha organizzato degli incontri mirati con le aziende proponenti dei progetti verso cui si sentiva maggiormente orientato. Questi confronti diretti hanno permesso di verificare la disponibilità dei referenti e risolvere le ultime incertezze, guidando il team verso una decisione finale pienamente informata e concorde.
\end{itemize}

\section{Analisi Capitolato selezionato}
\subsection{C5 - Nexum}

\section{Analisi altri Capitolati}
\subsection{C1 - Automated EN18031 Compliance Verification}

Il capitolato C1, proposto da Bluewind S.r.l., è stato inizialmente considerato dal gruppo per la sua aderenza a un tema attuale e rilevante: la conformità alla norma EN 18031, armonizzata rispetto alla direttiva RED per i dispositivi radio immessi nel mercato europeo.


\subsubsection{Obiettivo}

L'obiettivo del progetto è la realizzazione di un'applicazione in grado di supportare e automatizzare la verifica di conformità ai requisiti EN 18031 tramite esecuzione guidata di Decision Tree. Per ogni requisito analizzato, il sistema deve produrre un esito tra \textit{Pass}, \textit{Fail} e \textit{Not Applicable}, riducendo il carico di lavoro manuale e migliorando la tracciabilità delle decisioni.

In termini operativi, il software deve guidare l'utente nell'attraversamento dei nodi decisionali (domande con risposta \textit{Yes/No}) fino al nodo terminale e registrare in modo coerente percorso seguito, motivazioni e risultato finale per ciascun requisito.

\subsubsection{Caratteristiche funzionali principali}

Il capitolato richiede in particolare:

\begin{itemize}
    \item importazione di documenti tecnici e file esterni descrittivi dei Decision Tree (ad esempio XML e JSON);
    \item esecuzione interattiva degli alberi decisionali nel rispetto delle dipendenze gerarchiche tra requisiti;
    \item visualizzazione aggregata dei risultati tramite dashboard;
    \item possibilità di modifica e aggiornamento dei Decision Tree, con salvataggio in formati standard.
\end{itemize}

Oltre a queste funzionalità, il capitolato evidenzia anche l'esigenza di standardizzare la produzione della documentazione tecnica derivante dalla valutazione, così da rendere più semplice il riesame delle decisioni e il loro aggiornamento in caso di revisione normativa.

Il caso studio proposto (macchina del caffè connessa via Wi-Fi con comunicazione MQTT) rende il contesto concreto e utile per validare il comportamento dell'applicazione sui requisiti ACM e AUM. La presenza di un dominio applicativo reale costituisce un punto di forza, poiché permette di testare il flusso end-to-end su uno scenario non artificiale.

\subsubsection{Requisiti obbligatori e opzionali}

Dal punto di vista della pianificazione, il capitolato presenta una distinzione chiara tra requisiti obbligatori e opzionali.

Tra i requisiti obbligatori risultano centrali: importazione di formati standard, esecuzione automatizzata dei Decision Tree, rispetto delle dipendenze gerarchiche e dashboard di visualizzazione dei risultati.

Tra i requisiti opzionali emergono invece funzionalità di maggiore valore aggiunto (editor integrato, esportazioni in più formati, giustificazione dei risultati, editor grafico avanzato), che possono incrementare qualità del prodotto ma anche complessità implementativa.

Questa separazione è stata valutata positivamente dal gruppo perché facilita la definizione di una baseline minima e consente un'estensione incrementale del progetto.

\subsubsection{Considerazioni tecnologiche}

Il proponente non impone vincoli rigidi su stack, architettura o tipologia applicativa (desktop o web), lasciando libertà progettuale al team. L'unica preferenza esplicita riguarda l'uso di Python 3.x lato backend, preferibilmente con gestione del progetto tramite strumenti di packaging moderni.

Questa flessibilità rappresenta contemporaneamente un vantaggio e un costo: da un lato permette al gruppo di scegliere soluzioni aderenti alle proprie competenze, dall'altro richiede una fase preliminare più ampia di analisi e confronto tecnico prima di convergere su un'architettura definitiva.

\subsubsection{Incontro con il proponente}

\subsubsection{Valutazione del gruppo}

Durante la valutazione, il gruppo ha riconosciuto i punti di forza del capitolato: requisiti ben strutturati, dominio applicativo realistico e disponibilità del proponente a fornire supporto continuo sia da remoto sia in presenza.

Sono stati inoltre considerati positivamente:

\begin{itemize}
    \item la chiarezza della traccia progettuale e dei criteri di completamento;
    \item la presenza di un caso studio concreto per validare il prodotto;
    \item la forte rilevanza industriale e normativa del problema affrontato.
\end{itemize}

Parallelamente, il gruppo ha individuato alcune criticità:

\begin{itemize}
    \item curva di apprendimento iniziale elevata legata all'interpretazione della norma;
    \item forte dipendenza dalla corretta modellazione logica dei Decision Tree e delle loro gerarchie;
    \item margine percepito più limitato, rispetto ad altri capitolati, in termini di sperimentazione creativa e ampiezza delle funzionalità innovative.
\end{itemize}

Nonostante ciò, il progetto non è stato selezionato come scelta finale. Le principali motivazioni sono state:

\begin{itemize}
    \item interesse complessivamente più limitato del team verso il dominio normativo rispetto ad altri capitolati;
    \item percezione di un minor margine creativo e di estensione funzionale rispetto alle alternative considerate;
    \item necessità di un investimento iniziale significativo nell'interpretazione della normativa, ritenuto meno allineato agli obiettivi formativi prioritari del gruppo.
\end{itemize}

In conclusione, il capitolato C1 è stato valutato come solido, ben definito e supportato da un proponente disponibile, ma complessivamente meno allineato alle preferenze tecniche e agli interessi progettuali del gruppo aLittleByte rispetto al capitolato poi selezionato.

\subsection{C2 - Code Guardian}

Il capitolato C2, proposto da Var Group S.p.A., presenta lo sviluppo di una piattaforma web orientata all'audit dei repository software, con particolare attenzione a qualità del codice, sicurezza e manutenibilità.

\subsubsection{Obiettivo}

L'obiettivo del progetto è realizzare un sistema modulare, basato su un'architettura ad agenti coordinati da un orchestratore centrale, in grado di analizzare repository GitHub e produrre report automatici sullo stato complessivo del progetto.

La piattaforma deve non solo individuare criticità tecniche e di sicurezza, ma anche suggerire azioni correttive (\textit{remediation}) per migliorare progressivamente qualità, robustezza e mantenibilità del software analizzato.

\subsubsection{Caratteristiche funzionali principali}

Dal punto di vista funzionale, il capitolato richiede in particolare:

\begin{itemize}
    \item analisi automatica dei repository per identificare stack tecnologico, dipendenze e versioni utilizzate;
    \item verifica della presenza e della qualità dei test, con attenzione alla copertura e all'affidabilità complessiva della suite;
    \item audit di sicurezza orientato alle linee guida OWASP, con evidenziazione di vulnerabilità e configurazioni rischiose;
    \item valutazione della qualità della documentazione tecnica (README, API docs, guide operative);
    \item dashboard web per visualizzare stato generale, criticità e priorità di intervento.
\end{itemize}

Un elemento distintivo del capitolato è la componente di remediation: il sistema non si limita a segnalare problemi, ma propone interventi migliorativi concreti, favorendo un approccio operativo e continuo alla qualità del codice.

\subsubsection{Requisiti obbligatori e opzionali}

La documentazione propone un percorso strutturato che include attività iniziali di analisi e progettazione (raccolta requisiti, user story mapping, modellazione UML e schema dati), seguite dalla realizzazione di un MVP funzionante con demo e documentazione tecnica.

Dal punto di vista del gruppo, questa impostazione è stata considerata positiva perché orientata a un processo realistico di sviluppo, ma al tempo stesso impegnativa in termini di organizzazione e carico progettuale.

Tra le estensioni possibili, risultano particolarmente interessanti l'integrazione CI/CD più avanzata, l'analisi storica delle versioni, sistemi di ranking tra progetti e meccanismi interattivi di notifica e remediation.

\subsubsection{Considerazioni tecnologiche}

Il capitolato indica con chiarezza lo stack di riferimento: Node.js e Python per backend e orchestrazione, React.js per il frontend, MongoDB/PostgreSQL per la persistenza, GitHub Actions per l'automazione e AWS per l'infrastruttura cloud.

La chiarezza tecnologica è stata valutata favorevolmente perché riduce l'ambiguità iniziale e consente una pianificazione più precisa. Tuttavia, l'eterogeneità dello stack richiede competenze distribuite su più aree (sviluppo web, automazione, cloud, sicurezza), aumentando la complessità complessiva del progetto.

\subsubsection{Valutazione del gruppo}

Il gruppo ha riconosciuto in C2 diversi punti di forza:

\begin{itemize}
    \item tema attuale e di forte utilità pratica nel contesto industriale;
    \item architettura moderna e potenzialmente molto formativa;
    \item obiettivi ben definiti e tecnologie chiaramente esplicitate.
\end{itemize}

Parallelamente, sono emerse alcune criticità rilevanti:

\begin{itemize}
    \item complessità tecnica elevata dovuta alla combinazione di multi-agente, sicurezza applicativa e integrazioni esterne;
    \item carico consistente legato ai requisiti di qualità e validazione del software;
    \item interesse complessivo del gruppo inferiore rispetto ad altri capitolati ritenuti più allineati alle preferenze progettuali condivise.
\end{itemize}

In conclusione, C2 è stato considerato un capitolato solido, moderno e professionalmente molto valido, ma non è stato scelto dal gruppo aLittleByte poiché valutato meno aderente all'equilibrio desiderato tra interesse, complessità e obiettivi formativi.

\subsection{C4 - L'app che Protegge e Trasforma}
\subsection{C6 - Second Brain}

\end{document}
